\begin{table}[!ht]
\centering
\caption{Balanceamento entre trabalhadores brancos e negros \emph{antes} de qualquer
controle --- a tabela que mostra de que é feita a comparação (Angrist \& Pischke, cap.~3).
Diferença padronizada: $d = (\bar x_b - \bar x_n)/\sqrt{(s_b^2+s_n^2)/2}$; valores acima de
$0{,}10$ em módulo indicam desequilíbrio relevante. População completa da PEA com renda
positiva ($N = 7.699.278$; 3.271.443 brancos e 4.427.835 negros).}
\label{tab:balanceamento}
\small
\begin{tabular}{lrrrr}
\toprule
Variável & Brancos & Negros & Diferença & $d$ \\
\midrule
log do rendimento & 7,872 & 7,447 & 0,425 & 0,49 \\
Idade (centrada) & −1,637 & −3,120 & 1,483 & 0,11 \\
Sexo feminino & 0,437 & 0,399 & 0,038 & 0,08 \\
Fundamental completo ou mais & 0,790 & 0,687 & 0,103 & 0,24 \\
Médio completo ou mais & 0,660 & 0,531 & 0,129 & 0,27 \\
Superior completo ou mais & 0,276 & 0,142 & 0,134 & 0,34 \\
Pós-graduação & 0,083 & 0,036 & 0,047 & 0,20 \\
log(horas trabalhadas) & 3,645 & 3,601 & 0,044 & 0,11 \\
Reside em área urbana & 0,814 & 0,774 & 0,039 & 0,10 \\
Emprego formal & 0,494 & 0,441 & 0,053 & 0,11 \\
Conta própria & 0,276 & 0,281 & −0,005 & −0,01 \\
Trabalho doméstico & 0,049 & 0,079 & −0,030 & −0,12 \\
CBO: dirigente & 0,055 & 0,023 & 0,032 & 0,16 \\
CBO: profissional & 0,151 & 0,079 & 0,072 & 0,23 \\
CBO: técnico & 0,090 & 0,067 & 0,023 & 0,09 \\
\% negro na UPA ($z$) & −0,705 & 0,329 & −1,034 & −1,17 \\
Desemprego na UPA ($z$) & −0,331 & 0,091 & −0,422 & −0,46 \\
Educ. média na UPA ($z$) & 0,390 & −0,024 & 0,414 & 0,42 \\
\bottomrule
\end{tabular}
\normalsize
\par\smallskip
\footnotesize\emph{Suporte comum:} 97,2\% das UPAs têm
trabalhadores dos dois grupos, e 100,0\% das células
UF~$\times$~escolaridade também (100,0\% das observações
estão em células com os dois grupos) --- a comparação não depende de extrapolar para
regiões da amostra onde só existe um dos grupos.
\end{table}
